% ID: SEL-D0079-OB-E01
% ESTADO: cotejada. Segunda pasada visual y borrador compilado revisados.
% Fuente DOC: Selectividad/sel_2012_matematicas_aplicadas/1 MATEMÁT. APLIC. CC. SS/EXAMEN 1-ANDALUCÍA.doc
% SHA-256 original: a5dfff58dac253bf021b0d2b9409d68987e9ce6e97cfa4abbaf840019a7a18e2
% Página convertida: 2; opción B; ejercicio 1.

Sea el recinto determinado por las siguientes inecuaciones:
\[
3x+4y\geq28;\quad5x+2y\leq42;\quad x-y\geq0.
\]
\begin{enumerate}[label=\alph*)]
\item \textbf{(0.5 puntos)} Razone si el punto de coordenadas $(7,3)$ pertenece al recinto.
\item \textbf{(1.5 puntos)} Represente dicho recinto y halle sus vértices.
\item \textbf{(0.5 puntos)} Calcule el valor máximo de la función $F(x,y)=3x-2y+6$ en el recinto, indicando el punto o puntos donde se alcanza ese máximo.
\end{enumerate}
